 \documentclass[11pt]{article}
\usepackage{amsmath,amsfonts}
\usepackage{graphicx}
%\usepackage{xcolor}
%\definecolor{gray}{rgb}{.75,.75,.75}
%\usepackage[landscape]{geometry}
%\usepackage{hyperref}
%\usepackage[bookmarks=true, pdfborder={0 0 .5}, urlbordercolor=gray]{hyperref}

\renewcommand{\thefootnote}{\fnsymbol{footnote}}


\addtolength{\topmargin}{-1.0in}
\addtolength{\textheight}{2.0in}
\addtolength{\oddsidemargin}{-.5in}
\addtolength{\textwidth}{1in}
\pagestyle{empty}

\newcommand{\ds}{\displaystyle}
\newcommand{\ts}{\textstyle}

\newcommand{\Z}{\mathbb{Z}}
\newcommand{\R}{\mathbb{R}}
\newcommand{\C}{\mathbb{C}}
\newcommand{\EE}{\mathcal{E}}
\newcommand{\tZ}{\tilde{\Z}}

\newcommand{\Sym}{\mathrm{Sym}}

\renewcommand{\circle}[1]{{\bigcirc\hspace{-.75em}#1}}

\begin{document}

\footnote[0]{Notes by Peter Magyar \ {\tt magyar@math.msu.edu}}
\vspace{-1em}

\centerline{\bf Math 133\hspace{1.5em} \hfill{\Large\bf
Polar Areas and Lengths}\hfill 
Stewart \S10.4}
\vspace{1em}

\noindent  {\bf Slope in polar coordinates.}
We have seen that round, turny shapes are more simply described by polar $r\theta$-equations
than by rectangular $xy$-equations. In this section, we use 
polar equations to compute geometric information.

Thus, we consider a polar curve $r=f(\theta)$ over $\theta\in[a,b]$.
We split the interval $\theta\in[a,b]$ into a large number $n$ of increments,
each of length $\Delta\theta=\frac{b-a}n$,
with sample points $\theta_1,\ldots,\theta_n$.
Here is a typical increment of the curve over $\theta\in[\theta_i, \theta_{i{+}1}]$, 
showing the corresponding increments in the coordinates:
\\[1em]
\centerline{
\includegraphics[width=2.5in]
{10-4-Delta-Y.png} 
}
\vspace{-1em}

Our first problem is to find the slope of this curve at a given $\theta$.
It is {\it not} the derivative $f'(\theta)=\frac{dr}{d\theta}$, 
which is the rate of change of the radius with respect to the angle.
Rather, slope is the rate of change of $y = r\,\sin(\theta) = f(\theta)\,\sin(\theta)$
with respect to $x = r\,\cos(\theta) = f(\theta)\,\cos(\theta)$. That is:
$$
\text{(slope at $\theta$)} 
\ =\
\frac{dy}{dx}
\ =\ 
\frac{\ \ \frac{dy}{d\theta}\ \ }{ \frac{dx}{d\theta} }
\ =\
\frac{(f(\theta)\,\sin(\theta))'}{( f(\theta)\,\cos(\theta))' }
\ =\ \frac{f'(\theta)\,\sin(\theta) + f(\theta)\,\cos(\theta)}
{f'(\theta)\,\cos(\theta) - f(\theta)\,\sin(\theta)}\,.
$$
\\[0em]
{\bf Area in polar coordinates.} Assume $r=f(\theta)\geq 0$ for $\theta\in[a,b]$ to avoid complications with signs, and consider the region inside the curve, defined by $0\leq r\leq f(\theta)$ for $\theta\in[a,b]$.
Apply Slice Analysis (\S5.2), splitting the area $A$ into  
$n$ thin wedges $\Delta A_i$ over $[\theta_i,\theta_{i+1}]$:
\\[1em]
\centerline{
\includegraphics[width=2.5in]
{10-4-Delta-A.png} 
}
\vspace{-1em}

\noindent
We must compute the wedge area $\Delta A_i$.
Since $\Delta \theta$ is tiny, the small curve segments 
are very close to straight lines, and $\Delta A_i$ is a very thin triangle.
Neglecting the small piece with radius larger that $r_i$,
the slice $\Delta A_i$ is approximately an isosceles triangle with height $r_i$
and base $r_i\Delta\theta$.\footnote{
On a circle of radius $r$, and arc of $\theta$ radians has length $r\theta$:
this is the definition of radian measure.
}
Thus:
$$
\Delta A_i\ 
\approx\ 
\tfrac12\text{(base)}{\times}\text{(height)}
\ \approx\ 
\tfrac12(r_i\Delta\theta)r_i
\ =\ 
\tfrac12 r_i^2\Delta\theta
%\ =\
%\tfrac12 f(\theta_i)^2\Delta\theta
\,.
$$
The total area is the sum of these pieces, which is clearly a Riemann sum for an integral:
$$
A \ =\ \lim_{n\to\infty} \sum_{i=1}^n \Delta A_i
\ =\ \lim_{n\to\infty} \sum_{i=1}^n \tfrac12 r_i^2\Delta\theta
\ =\ \lim_{n\to\infty} \sum_{i=1}^n \tfrac12 f(\theta_i)^2\Delta\theta
%\ =\ \int_a^b \tfrac12 r^2\,d\theta
\ =\ \int_a^b \tfrac12 f(\theta)^2\,d\theta\,.
$$
That is, the area inside a polar graph $r=f(\theta)$ is given by an integral formula, but a different integral 
from the area under a rectangular graph $y=f(x)$.
\\[1em]
{\bf Arclength in polar coordinates.} 
Finally, we compute the length of the curve $r=f(\theta)$ for $\theta\in[a,b]$.
The length $L$ is a sum of $n$ increments $\Delta L_i$:
\\[1em]
\centerline{
\includegraphics[width=3in]
{10-4-Delta-L.png} 
}
\vspace{-1em}

\noindent
Each increment $\Delta L_i$  is approximately a straight line segment.
Next to it is the radial segment $\Delta r$ and
the tiny circular arc with length $r_i\,\Delta\theta$, 
which is also approximately a straight line.
We get an approximate right triangle with hypotentuse $\Delta L_i$
and legs $r_i\,\Delta\theta$ and $\Delta r$, so the Pythagorean Theorem gives:
$$
\Delta L_i \ 
\approx\ \sqrt{(r_i\Delta \theta)^2+(\Delta r)^2}
\ =\ 
\sqrt{
\tfrac {(r_i\Delta \theta)^2+(\Delta r)^2} {(\Delta \theta)^2}
}\,\Delta\theta
\ =\ 
\sqrt{ 
 r_i^2 + (\tfrac{\Delta r}{\Delta \theta})^2
} 
\,\Delta\theta\,.
$$
Therefore the total arclength is:
$$
L \ =\ \lim_{n\to\infty} \sum_{i=1}^n \Delta L_i
\ =\ \lim_{n\to\infty} \sum_{i=1}^n 
\sqrt{ r_i^2 + (\tfrac{\Delta r}{\Delta\theta})^2 }\,\Delta\theta
$$
$$
\ =\ \lim_{n\to\infty} \sum_{i=1}^n 
\sqrt{ f(\theta_i)^2 + (\tfrac{\Delta f(\theta_i)}{\Delta\theta})^2 }\,\Delta\theta
%\ =\ 
%\int_a^b \sqrt{ r^2 + (\tfrac{dr}{d\theta})^2 }\,d\theta
\ =\ 
\int_a^b \sqrt{f(\theta)^2 + f'(\theta)^2}\,d\theta\,.
$$
We could also deduce this from our previous parametric arclength formula (\S10.3) by applying
it to $(x(t),y(t))=(f(t)\cos(t), f(t)\sin(t))$.

\pagebreak

\noindent
{\bf Example: Area of intersections.} Consider the polar curve $r = f(\theta)=1-\cos(\theta)$.
We picture the abstract function $f$ by its rectangular graph in $\theta r$ parameter space 
(end \S10.3):
\\[1em]
\centerline{
\includegraphics[width=2.7in]
{10-4-Cardioid-Rect.png} 
}
\vspace{-.8em}

\noindent
The polar graph is a cardioid (heart-shape), which we draw along with the circle $r=\frac12$.
\\[.5em]
\centerline{
\includegraphics[width=3in]
{10-4-Cardioid.png} 
}
\vspace{-.8em}

\noindent
{\sc prob:} Find the area of the crescent-shaped region
inside the cardioid \& outside the circle.
\\[.5em]
We must first determine the intersection points of the two curves, where:
$$
r = 1-\cos(\theta) = \tfrac12
\quad\Longrightarrow\quad
\cos(\theta)=\tfrac12
\quad\Longrightarrow\quad
\theta= \pm\tfrac\pi3 + 2n\pi\,,
$$
where $n$ is any integer. Since the whole cardioid is traced by $\theta\in[0,2\pi]$,
we can take all intersection points in this range: 
$\theta=\frac\pi3$ and $\theta = -\frac{\pi}3 + 2\pi = \frac{5\pi}{3}$.
Now we take the area inside the cardioid $r=f(\theta)=1-\cos(\theta)$, 
minus the area inside the circle $r=g(\theta)=\frac12$:
$$
A 
\ =\ \int_a^b \tfrac12 f(\theta)^2 -\tfrac12 g(\theta)^2\ d\theta
\ =\ \int_{\pi/3}^{5\pi/3} \tfrac12(1{-}\cos(\theta))^2 - \tfrac12(\tfrac12)^2\ d\theta
$$
$$
\ =\
\left[\vphantom{\sum}\tfrac58\theta- \sin(\theta) + \tfrac18\sin(2\theta)\,\right]_{\theta=\pi/3}^{\theta=5\pi/3}
\ =\
\tfrac78\sqrt 3 +\tfrac{5\pi}6
\ \approx\ 4.1\,.
$$
To do the integral, expand $(1-\cos(\theta))^2$ and use
$\cos^2(\theta) = \frac12 + \frac12\cos(2\theta)$ (see \S7.2).
\\[1em]
{\sc prob:} Find the highest point of the cardioid: $y(t) = (1-\cos(t))\sin(t) = \text{max}$. From \S3.1, we need to solve
$y'(t) =0$. Using the Product Rule and $\sin^2=1-\cos^2$ we get:
 $y'(t)=\sin(t)\sin(t) + (1-\cos(t))\cos(t) = -2\cos^2(t)+\cos(t)+1=0$, and the Quadratic formula gives
 $\cos(t) = -\frac12$, $t=\pm \frac{2\pi}{3}$. 
The maximum is $y(\frac{2\pi}{3}) = \frac{3\sqrt 3}{4}$.

\pagebreak

\noindent
{\bf Review example: Exponential spiral.} Consider a snail-shell spiral curve which
doubles in radius with each turn:
\\[0em]
\centerline{
\includegraphics[width=2.5in]
{10-4-Double-Spiral.png} 
}
\vspace{-.5em}

\noindent
This is the polar graph  $r = f(\theta)=ca^\theta=ce^{b\theta}$
of a general exponential function (\S6.4). 
Assuming $f(0)=1$, $f(2\pi)=2$, allows us to solve 
for $c=1$ and $a= 2^{1/2\pi} =e^{\ln(2)/2\pi}$ to get:
$$
r\,=\,2^{\,\theta/2\pi} \,=\, e^{b\theta}
\quad\text{for}\quad b=\tfrac{\ln(2)}{2\pi}\,.
$$
What is the length of this curve, from the point $(r,\theta)=(1,0)$ all the way to the center,
that is, for $\theta\in (-\infty,0]$? 
We have $f'(\theta)  \,=\, (e^{b\theta})' = be^{b\theta}$, so
the arclength formula gives:
$$
L 
\ = \ 
\int_{-\infty}^0\!\!\!\! \sqrt{f(\theta)^2 + f'(\theta)^2}\,d\theta
\ = \ 
\int_{-\infty}^0\!\!\!\! \sqrt{1{+}b^2}\, e^{b\theta}\,d\theta
\ = \ 
\left[ \tfrac1b\sqrt{1{+}b^2}\, e^{b\theta}\right]_{\theta=-\infty}^{\theta=0}
$$
$$
\ = \ 
\tfrac1b\sqrt{1{+}b^2} \,-\, \lim_{N\to\infty} \tfrac1b\sqrt{1{+}b^2}\, e^{-N}
\ =\
\sqrt{1{+}\tfrac{4\pi^2}{\ln^2(2)}} \ \approx\ 9.12\,.
$$
Or we could use geometry to show that these infinitely many turns have finite length. 
Let $L_1$ be the length of the first turn $\theta\in[-2\pi,0]$, and $L_2$ the length of the second turn, etc. 
The exponential spiral is {\it scale invariant}: 
each turn inward is the $\frac12$ dilation of the previous turn,
with half the length, so the total is a geomteric series 
$\sum_{n=1}^\infty cr^{n-1} = \frac c{1-r}$ (\S11.2):
$$
L \,=\, L_1 +L_2+L_3+L_4+\cdots
\,=\,  L_1+\frac12 L_1 +\frac1{2^2}L_1+\frac1{2^3}L_1+\cdots
\,=\, \frac{L_1}{(1-\frac12)} \,=\, 2L_1.
$$
{\bf Harmonic Spiral.} From the above, we may say that the inward spiral $r=1/2^{\theta}$ 
 has finite arclength as $\theta\to\infty$ because the geometric series 
 $\sum_{n=1}^\infty \frac1{2^n}$ is  convergent.
Let us instead model an inward spiral on the divergent harmonic series $\sum\frac1n=\infty$, 
namely $r=\frac1{\theta}$ for $\theta\geq 1$. Then this should have infinite arclength:
$$\textstyle
L \ =\ \int_1^{\infty} \sqrt{(\frac{1}{\theta})^2+(-\frac{1}{\theta^2})^2}\,d\theta
\ =\ \int_1^{\infty} \frac{1}{\theta}\sqrt{1+\frac{1}{\theta^2}}\,d\theta
\ =\ \int_1^{\infty} \frac{\sqrt{1+\theta^2}}{\theta^2}\,d\theta.
$$
Since the integrand (in the second integral) is clearly positive and decreasing, the Integral Test  (\S11.3) 
tells us that this diverges whenever the corresponding series diverges, namely :
$$\textstyle
\sum_{n=1}^\infty  \frac{1}{n}\sqrt{1+\frac{1}{n^2}}
\  \approx\  \sum_{n=1}^\infty  \frac{1}{n} \ = \ \infty\  \text{(divergent)}.
$$
This can be justified by the Direct Comparison Test (\S11.4), since 
$\tfrac{1}{n}\sqrt{1{+}\tfrac{1}{n^2}}>\tfrac{1}{n}$; or 
the Limit Comparison Test: the ratio $\frac{a_n}{b_n}=\sqrt{1{+}\tfrac{1}{n^2}}\to 1$, 
so they have the same divergence.

\newpage

\noindent
Alternatively, we can directly integrate, 
switching the variable to $\int\! \frac{\sqrt{1+x^2}}{x^2}\,dx$. 
Since we have $\sqrt{1{+}x^2}$ (\S7.3), we try the trig substitution 
$x = \tan(t)$, $\sqrt{1{+}x^2}=\sec(t)$, $dx = \sec^2(t)\,dt$:
$$
\int\! \frac{\sqrt{1+x^2}}{x^2}\,dx
\ =\ 
\int\! \frac{\sec(t)}{\tan^2(t)}\sec^2(t)\,dt
\ = \ 
\int\! \frac{1}{\sin^2(t)\cos(t)}\,dt
\ = \ 
\int\! \frac{1}{\sin^2(t)\cos^2(t)}\cos(t)\,dt.
$$
As in \S7.2, we do the substitution $u=\sin(t)$, $1{-}u^2 = \cos^2(t)$, $du = \cos(t)\,dt$:
$$
\int\! \frac{1}{\sin^2(t)\cos^2(t)}\cos(t)\,dt
\ =\ 
\int\! \frac{1}{u^2(1{-}u^2)}\ du 
\ =\ \int\! \frac{A}{u^2}+\frac{B}{u}+\frac{C}{1{+}u}+\frac{D}{1{-}u}\ du.
$$
Here the result is a rational function, expanded by partial fracitions (\S7.4).
Then: 
$$
\frac{1}{u^2(1{-}u^2)}
\ =\ 
 \frac{1{-}u^2}{u^2(1{-}u^2)} + \frac{u^2}{u^2(1{-}u^2)}
 \ =\ \frac{1}{u^2} + \frac{1}{(1{-}u)(1{+}u)}
  \ =\ \frac{1}{u^2} + \frac{C}{1{+}u}+\frac{D}{1{-}u}.
$$
We can find the remaining coefficients by clearing denominators to get 
$$1= C(1{-}u)+D(1{+}u).$$
Substituting $u = 1$ gives $D=\frac12$, and $u=-1$ gives $C=\frac12$.
The integral becomes:
$$
\int \frac{1}{u^2}+\frac{\frac12}{1{+}u}+\frac{\frac12}{1{-}u}\ du
\ =\ 
-\frac{1}{u} + \tfrac12 \ln(1{+}u) - \tfrac12 \ln(1{-}u)
\ =\ 
-\frac{1}{u} + \frac12 \ln\!\left(\frac{1{+}u}{1{-}u}\right).
$$
Now we need to restore the original variable $x=\tan(t)$ from $u=\sin(t)$.
The standard triangle for $x=\tan(t)$ implies $u=\sin(t) = \frac{x}{\sqrt{1+x^2}}$.
After simplification, the final answer is:
$$
\int\! \frac {\sqrt{1{+}x^2}} {x^2}\,dx
\ =\ -\frac{\sqrt{1{+}x^2}}x 
+ \frac12\ln\!\left(\!\frac {\sqrt{1{+}x^2}+x} {\sqrt{1{+}x^2}-x} \!\right).
$$
Therefore the total arclength is:
$$
L \ =\ 
\int_1^\infty\! \frac {\sqrt{1{+}x^2}} {x^2}\,dx
\ =\ 
\lim_{x\to\infty} \tfrac12   \ln\!\left(\!\frac {\sqrt{1{+}x^2}+x} {\sqrt{1{+}x^2}-x} \!\right)
+ K
$$
for a constant $K$.  In the fraction $\frac {\sqrt{1{+}x^2}+x} {\sqrt{1{+}x^2}-x}$, 
we {\it cannot} use L'Hopital's Rule (\S6.8), 
since the numerator clearly approaches $\infty$,
but the denominator does not. Substitute $x=1/z$:
$$
\lim_{x\to\infty} \sqrt{1{+}x^2}-x \ =\ 
\lim_{z\to0^+}\! \sqrt{1{+}\frac1{z^2}}-\frac 1z \ =\ 
\lim_{z\to0^+}\! \frac{\sqrt{z^2{+}1}-1}z\,. 
$$
This is a $\frac 00$ limit, so we {\it can} apply L'Hopital to get:
$$
\lim_{x\to\infty} \sqrt{1{+}x^2}-x 
\ =\ \lim_{z\to0^+}\! \frac{(\sqrt{z^2{+}1}-1)'}{(z)'}
\ =\ \lim\limits_{z\to0^+}\! \frac {\frac z{\sqrt{z^2+1}}}   {1} \ =\ 0^+. 
$$
After all that, we obtain the expected arclength:
$$
L \ =\ 
\tfrac12   \ln\!\left(\!\frac {\infty} {\ \,0^+} \!\right) + K \ =\ 
\tfrac12   \ln(\infty) + K \ =\ \infty\,.
$$


\end{document}


$$\textstyle
L \ =\ \int_1^{\infty} \sqrt{(\frac{1}{\theta})^2+(-\frac{1}{\theta^2})^2}\,d\theta
\ =\ \int_1^{\infty} \frac{1}{\theta}\sqrt{1+\frac{1}{\theta^2}}\,d\theta.
$$

Or we could directly integrate, substituting
$z = \sqrt{1{+}\frac{1}{\theta^2}}$, 
$\theta = \frac{1}{\sqrt{z^2{-}1}}$,
$d\theta = -\frac{z}{(z^2-1)^{3/2}} dz$:
$$\textstyle
\int \frac{1}{\theta}\sqrt{1{+}\frac{1}{\theta^2}}\,d\theta
= -\int\! \sqrt{z^2{-}1}\cdot z\cdot \frac{z}{(z^2-1)^{3/2}}\,dz
%= -\int\! \frac{z^2}{z^2-1}\,dz
\stackrel{\text{P.F.}}= \int\! -1{+}\frac{1/2}{z+1}{-}\frac{1/2}{z-1}\,dz
=-\sqrt{1{+}\frac{1}{\theta^2}}
\,+\,\frac12\ln\!\frac{\sqrt{1+\theta^2}+\theta}{\sqrt{1+\theta^2}-\theta}.
$$
To find $\lim\limits_{\theta\to\infty}\sqrt{1{+}\theta^2}-\theta$, 
substitute $\theta=\frac1t$, getting
$
\lim\limits_{t\to 0^+}\frac{\sqrt{t^2{+}1}-1}t
\stackrel{\text{Hop}}=
\lim\limits_{t\to 0^+}\frac t{\sqrt{t^2{+}1}} =0^+
$.
\\
Hence the integral is $L=-1-(-1-\frac12\ln\frac{\infty}{0^+})=+\infty$.


%%%%%%%%%%%  Picture  %%%%%%%%%%
\\[1em]
\centerline{
%\includegraphics[width=2in]
{1-8-Discont-iv.png}
\qquad \qquad 
%\includegraphics[width=2in]
{1-8-Discont-v.png} 
}
\vspace{0em}

\noindent
%%%%%%%%%%%%%%%%%%%%%%%%%

%%%%%%%%%%%  Table  %%%%%%%%%%
In fact, we get the following derivatives:
$$\begin{array}{|c||c|c|c|c|c|c|}
\hline &&&&&& \\[-.9em] 
f(x) & \sin(x) & \cos(x) & \tan(x)  & \sec(x) & \csc(x) & \cot(x)
\\[.2em] \hline &&&&&& \\[-.9em]
f'(x) &\cos(x) & -\sin(x) & \sec^2(x) & \tan(x)\sec(x) & -\cot(x)\csc(x) & -\csc^2(x)
\\[.1em] \hline
\end{array}$$
\\[1em]
%%%%%%%%%%%%%%%%%%%%%%%%%%



















